\documentclass{exam}[12pt]

\usepackage{amsfonts,amsmath,amsthm,amssymb}


\boxedpoints

\addpoints
\pointsinmargin

\begin{document}
\pagestyle{head}

\firstpageheader{\large Calculus II}{\bf\Large Worksheet}{\large Name:\hspace{1in} }
%\extraheadheight{.7in}
%\extrafootheight{-1in} 
\small{\textit{
%Calculators are not allowed on this  exam. 
%Work all questions completely.  Show all work as described in class. \\
Copyright 2006 Dr. Brock Williams and
Texas Tech University.  Unauthorized reproduction prohibited.
}}

\begin{questions}
\large{


\question If the ammount of mayhem is constant in an episode of the
Muppet Show, there is a simple formula that gives the amount of mess created:
\[
\text{Mess = Mayhem * Time}
\]
That is, the mess is the product of the mayhem and the length of time it lasts.
Typically, however, the mayhem varies over time.  
In a certain 30 minute episode the amount of mayhem was given by the function 
$M(t) = 4 \ln(t+1)$, where $0\leq t\leq 30$.
Describe in complete detail how integration can be used to 
find the total mess in the episode.
\vspace{.5in}

\question  It is a well-known fact that the anger generated by a fan
during a football game is caused by both the frustration they experience
and the length of time it occurs.  The anger produced by a 
 constant ammount of frustration is given by the formula
\[
\text{Anger = }\frac{1}{3} \left(\text {Frustration}\right)^2 * \text{Time}.
\]
In the first quarter of a game, the Frustration was given by
$\displaystyle F(t)= \sin\left(\frac{\pi}{30} t\right)  $, where $0\leq t\leq 15$.
\vspace{.2in}

\textbf{Describe in complete detail} how integration can be used to find
the total Anger in the first quarter.  Discuss the role played by Riemann 
sums, and set up an integral to find this Anger.

\vspace{.7in}

\question Set up integrals to find the mass, moment about the $x$-axis, and
moment about the $y$-axis of the lamina bounded by $y=\sqrt{x}$ and
$y=\displaystyle \frac{1}{3} x$.  
Give the formula to find the coordinates of the center of mass.
\vspace{.5in}

\question A cylindrical tank has a height of 9 feet and a radius of 5 feet.
It is filled with a fluid having density 95 lb/ft$^3$.
Set up an integral to find the work done in pumping the fluid out of the tank if
\begin{parts}
\part The tank is full.
\vspace{.5in}
\part The tank is half full  and the water must also be
 pumped out a 3 foot pipe sticking
out the top of the tank.
\vspace{.5in}
\end{parts}

\question A conical tank  has a radius of 5 ft and a height of 13 ft.
Set up an integral to determine how much work is done in emptying it if
\begin{parts} 
\part It is filled with milk (which weighs 64.5 $lb/ft^3$).
\vspace{.5in}
\part It is filled to a depth of 9 ft and the milk must be pumped out
a 2 ft pipe at the top of the tank.
\vspace{.5in}
\end{parts}

\question  A dam has a triangular gate with dimensions as shown below.
The top of the gate is 3 feet below the water level.  Set up an integral
to find the force on
the gate. (Recall the density of water is 62.4 lb/ft$^3$.)
\vspace{.5in}

\question %Discuss the following completely.
\begin{parts}
\part What is the fundamental idea of Chapter 6?  
\vspace{.5in}
\part Suppose an object moves along the $x$-axis from $x=2$ to $x=5$ and is
acted on by a force $F(x)=x^2$.  Set up an integral to determine the amount of
work done on the particle.
\vspace{.5in}
\part Use the fundamental idea of Chapter 6 to explain in detail why the work
on the particle is given by your integral. 
\end{parts}



}
\end{questions}
\end{document}



